\def\stat{zac+garanin}



\def\tit{ОЦЕНКА НАДЕЖНОСТИ ГИБРИДНОГО ВЫСОКОПРОИЗВОДИТЕЛЬНОГО ВЫЧИСЛИТЕЛЬНОГО 

КОМПЛЕКСА ПРИ~РЕШЕНИИ НАУЧНЫХ ЗАДАЧ$^*$\\[-5pt]}



\def\titkol{Оценка надежности ГВВК %гибридного высокопроизводительного вычислительного комплекса 

при~решении научных задач}





\def\aut{А.\,А.~Зацаринный$^1$, А.\,И.~Гаранин$^2$, В.\,А.~Кондрашев$^3$, 

К.\,И.~Волович$^4$, С.\,И.~Мальковский$^5$\\[-3pt]}



\def\autkol{А.\,А.~Зацаринный, А.\,И.~Гаранин, В.\,А.~Кондрашев и~др.} 

%К.\,И.~Волович$^4$, С.\,И.~Мальковский$^5$}



\titel{\tit}{\aut}{\autkol}{\titkol}



\index{Зацаринный А.\,А.}

\index{Гаранин А.\,И.}

\index{Кондрашев В.\,А.}

\index{Волович К.\,И.}

\index{Мальковский С.\,И.}

\index{Zatsarinny A.\,A.}

\index{Garanin A.\,I.}

\index{Kondrashev V.\,A.}

\index{Volovich K.\,I.}

\index{Malkovsky S.\,I.}





{\renewcommand{\thefootnote}{\fnsymbol{footnote}}

\footnotetext[1]

{Работа выполнена при частичной поддержке РФФИ (проекты 18-29-03100

и~18-29-03196).}}



\renewcommand{\thefootnote}{\arabic{footnote}}

\footnotetext[1]{Институт проблем информатики Федерального исследовательского центра <<Информатика и~управление>> 

Российской академии наук, \mbox{AZatsarinny@frccsc.ru}}

\footnotetext[2]{Институт проблем информатики Федерального исследовательского центра <<Информатика и~управление>> 

Российской академии наук, \mbox{algaranin@mail.ru}}

\footnotetext[3]{Институт проблем информатики Федерального исследовательского центра <<Информатика и~управление>> 

Российской академии наук, \mbox{VKondrashev@frccsc.ru}}

\footnotetext[4]{Институт проблем информатики Федерального исследовательского центра <<Информатика и~управление>> 

Российской академии наук, \mbox{KVolovich@frccsc.ru}}

\footnotetext[5]{Вычислительный центр Дальневосточного отделения Российской академии наук, 

\mbox{sergey.malkovsky@gmail.com}}





 \label{st\stat}



                  \thispagestyle{headings}

                  

  \vspace*{-16pt}

  

           

      

      \Abst{Обоснована необходимость применения гибридных решений при создании 

высокопроизводительных вычислительных систем. Дано краткое описание гибридного 

высокопроизводительного вычислительного комплекса (ГВВК) Федерального исследовательского

центра <<Информатика и~управ\-ле\-ние>> Российской академии наук

(ФИЦ ИУ РАН) 

и~представлена его укрупненная структурная схема. Предложен методический подход 

к~оценке надежности ГВВК, на основе которого проведены расчеты надежности 

выделенных функциональных подсистем. Отдельно проведена оценка надежности 

<<вычислительной инфраструктуры>> ГВВК (без периферийных элементов). Даны 

рекомендации по повышению надежности функциональных подсистем.}

      

      \KW{гибридный высокопроизводительный вычислительный комплекс; 

функциональные подсистемы; надежность; отказ; эквивалентная схема для расчета 

надежности}



\DOI{10.14357\!/\!08696527190212} 





\vskip 10pt plus 9pt minus 6pt





      \thispagestyle{headings}

      

      \vspace*{-20pt} 

      

\section{Введение}



\vspace*{-3pt}



      В настоящее время одним из наиболее перспективных направлений развития 

высокопроизводительных вычислительных систем стали гибридные решения~[1], 

предусматривающие в~дополнение к~центральным процессорам установку различных 

<<ускорителей вычислений>>, или сопроцессоров. При подключении к~центральному 

процессору по высокоскоростной шине такие сопроцессоры берут на себя большую часть 

вычислительной работы.

      

      Идея гибридизации вычислительных систем на аппаратном уровне не нова. Еще 

в~1980-х~гг.\ в~некоторых вычислительных системах в~дополнение к~центральным 

процессорам (CPU, central processing units) стали устанавливать сопроцессоры, которые брали на себя часть 

вычислительных задач. В~качестве примера можно привести сопроцессор Intel~8087, 

выпущенный в~1980~г. Он использовался вместе с~центральными процессорами 

Intel~8086 для ускорения расчетов с~плавающей точкой. Еще раньше применялись 

сопроцессоры, ускоряющие операции вво\-да--вы\-вода.

      

      Новый этап развития гибридных вычислительных систем начался в~середине  

2000-х~гг., когда для вычислительных целей стали использовать графические 

процессоры (GPU, graphic processing units)~[2]. Интерес к~гибридным схемам обусловлен высокой 

производительностью GPU, которая объясняется особенностями их архитектуры. Если 

современные CPU содержат несколько ядер (на большинстве современных систем от~2 

до~4), то графический процессор изначально создавался как многоядерная структура, 

в~которой число ядер измеряется сотнями. Разница в~архитектуре обусловливает и~разницу в~принципах работы. Если архитектура CPU изначально предполагала последовательную 

обработку информации, то GPU предназначался для обработки компьютерной графики, 

так как рассчитан на  

мас\-сив\-но-па\-рал\-лель\-ные вычисления~[2]. Графические процессоры уже достигли той точки развития, 

когда многие вычислительные задачи могут с~легкостью решаться на них, причем 

быстрее, чем на многоядерных системах классической архитектуры. 

      

      Каждая из рассматриваемых архитектур имеет свои достоинства. Цннтральный

      процессор лучше 

работает с~задачами, имеющими небольшую степень параллелизма (большое чис\-ло 

проверок условий, ветвления, необходимость частой синхронизации потоков и~т.\,д.). При 

отсутствии таких ограничений вместе с~необходимостью обработки больших объемов 

информации очевидное преимущество приобретают GPU.

      

      Впервые появившись в~рейтинге top500 самых мощных вычислительных сис\-тем 

в~2006~г., в~настоящий момент (на ноябрь 2018~г.)\ гибридные вычислительные сис\-те\-мы 

обеспечивают 41\% суммарной производительности указанного мирового рейтинга. При 

этом в~первой десятке мирового рейтинга число подобных систем составляет уже 70\% 

и~будет в~дальнейшем неизбежно расти. 



      Среди вычислительных платформ подобного класса можно выделить архитектуру 

POWER, на основе которой создан ряд специализированных вычислительных систем, 

включающих самые мощные в~настоящее время суперкомпьютерные установки в~мире~--- 

Sierra и~Summit с~пиковой производительностью более~125 и~200~Пфлопс 

соответственно. 





\vspace*{-9pt}

      

\section{Краткое описание гибридного высокопроизводительного 

вычислительного комплекса Федерального исследовательского центра 

<<Информатика и~управление>> 

Российской академии наук}



\vspace*{-3pt}



      В ФИЦ ИУ РАН в~рамках программы развития преду\-смот\-ре\-но 

формирование современной исследовательской инфраструктуры для повышения 

эффективности проводимых фундаментальных научных исследований, а также решения 

актуальных на\-уч\-но-прак\-ти\-че\-ских задач.

      

      Для предоставления вычислительных ресурсов заинтересованным научным 

коллективам в~ФИЦ ИУ РАН на базе центра обработки данных (ЦОД) ФИЦ ИУ РАН 

создан ГВВК, 

предоставляющий вычислительную инфраструктуру научным подразделениям РАН. 



 \begin{figure} %fig1

 \vspace*{-2.5pt}   % чтобы не свистел float по вертикали

\begin{center}

\hspace*{12pt}\mbox{%

 \epsfxsize=96.849mm 

 \epsfbox{zac-1.eps}

 }

\hspace*{10pt}

\end{center}



\vspace*{-506pt}  % регулировка по высоте для подписи



\hspace{336pt}\hbox{\rotcaption{\protect\small\mdseries       % сдвиг подписи по горизонтали

Структурная схема ГВВК}}\hspace*{-13pt} % чтобы не свистело

      \end{figure}



      





      

      Гибридный высокопроизводительный вычислительный комплекс  создан на основе серверов архитектуры POWER различных модификаций 

(серверы Power~9 укомплектованы четырьмя графическими ускорителями Tesla~V100, 

а~также имеются серверы, оснащенные только классическими процессорами CPU). 

      

      Средствами ЦОД было организовано сетевое взаимодействие компонентов ГВВК 

между собой, с~пользователями системы и~вспомогательными сервисами, 

предоставляемыми облачной инфраструктурой ЦОД ФИЦ ИУ РАН. 





      

      Структурная схема гибридного высокопроизводительного вычислительного 

комплекса представлена на рис.~1.



\section{Методический подход к~оценке надежности гибридного высокопроизводительного вычислительного комплекса}



      При решении вопросов, связанных с~оценкой надежности систем, аналогичных 

ГВВК, предлагается 

учитывать следующие особенности~[3]:

      \begin{itemize}

      \item ГВВК является многофункциональной системой, функции которой имеют 

существенно различную значимость и, соответственно, характеризуются разным уровнем 

требований к~надежности их выполнения;

\item в~ГВВК возможно возникновение некоторых исключительных (аварийных, 

критических) ситуаций, представляющих сочетание отказов либо ошибок 

функционирования комплекса и~способных привести к~значительным нарушениям 

процессов функционирования (авариям);

      \item в~функционировании ГВВК участвуют различные виды ее обеспечения 

и~персонал ГВВК, которые могут в~той или иной степени влиять на уровень надежности 

ГВВК;

      \item  на процесс функционирования ГВВК оказывают влияние большое число 

разнородных факторов (технические средства, программное обеспечение, человеческий 

фактор и~др.), при этом в~выполнении одной функции ГВВК обычно участвуют несколько 

различных элементов, а один и~тот же элемент может участвовать в~выполнении 

нескольких функций системы.

      \end{itemize}

      

      С учетом этих особенностей при решении вопросов надежности ГВВК ее 

количественное описание, анализ, оценка и~обеспечение надежности предлагается 

проводить по каждой функции ГВВК в~отдельности~[3]. 

      

      Совокупность технических и~программных элементов ГВВК, выделяемая из всего 

состава ГВВК по признаку участия в~выполнении некоторой ($i$-й) функции системы, 

образует $i$-ю функциональную подсистему (ФП) ГВВК.

      

      Анализ надежности ГВВК в~реализации его функций проводится по каждой ФП 

ГВВК в~отдельности с~учетом уровня надежности и~других свойств входящих в~нее 

технических и~программных средств, а также человеческого фактора.

      

      Выбор состава показателей надежности ГВВК осуществляется на основе 

установленного перечня функций системы. Расчет надежности технических систем по 

безотказности обычно проводится в~предположении, что вся система и~каждый ее элемент 

могут находиться только в~одном из двух возможных состояний~--- работоспособном 

и~неработоспособном, а отказы элементов независимы друг от друга.

      

      Проведем структурно-логический анализ~[4, 5] процессов функционирования 

и~взаимосвязей элементов в~процессе функционирования ГВВК, 

структурная схема которого  представлена на рис.~1.

      



      

      Доступ удаленных пользователей осуществляется через сеть Интернет 

с~использованием средств обеспечения информационной безопасности (ИБ) ЦОД. 

Локальные пользователи подключаются через <<Средства доступа>> (коммутаторы~7 

и~8). Средства ИБ и~средства доступа подключаются к~коммутаторам~1 или~2. 

      

      При оценке надежности ГВВК выделим три ФП: 

      \begin{enumerate}[(1)]

      \item ФП для решения <<информационных>> задач;

      \item ФП для решения задач <<подготовки данных>>;

      \item  ФП для решения <<расчетных>> задач.

      \end{enumerate}

      

      Серверы 1 и~3 совместно с~системами хранения данных (СХД)~1 и~2 используются для 

решения задач <<информационного>> типа. При решении задач <<подготовки данных>> 

указанные серверы взаимодействуют только с~СХД~2.

      

      Серверы 4 и~5 под управлением сервера~2 во взаимодействии с~СХД~1 образуют высокопроизводительный вычислительный комплекс и~используются 

для решения расчетных задач.

      

      Коммутаторы 1 и~2 относятся к~локальной вычислительной сети (LAN, Local Area Network), 

коммутаторы~3--6~--- к~сети хранения данных (SAN, Storage Area Network).

      

      Принято, что если в~составе сервера имеется неисправный компонент, то такой 

сервер выводится из эксплуатации для ремонта.

      

      Сформулируем понятие <<отказ>>. Под \textbf{отказом} 

ГВВК будем понимать событие, 

заключающееся в~прекращении решения задач <<информационного>> типа, задач 

<<подготовки данных>> и~расчетных задач по причине неисправности  

ап\-па\-рат\-но-про\-грам\-мных средств и~требующее для восстановления процесса 

функционирования проведения ремонтных работ с~привлечением обслуживающего 

персонала. 

      

      Используя результаты структурно-логического анализа структурной схемы, 

представленной на рис.~1, и~определение понятия <<отказа>> системы, составим 

эквивалентные схемы для расчета надежности указанных выше ФП.

      

      Под \textit{эквивалентной схемой для расчета надежности} (<<надежностной 

структурой>>, <<структурной схемой надежности>>) системы понимается логическое 

соединение элементов в~таком виде, из которого следовало бы определение отказа 

системы~[4, 5]. 

      

      При построении эквивалентных схем будем исходить из расчета на <<наихудший 

случай>>, т.\,е.\ выбираем такой вариант построения, при котором для решения 

пользователем конкретной задачи задействуется наибольшее число элементов сис\-темы.

      

      Эквивалентная схема для расчета надежности ГВВК при решении 

<<информационных>> задач изображена на рис.~2, где под СКС понимается 

структурированная кабельная система, обеспечивающая соединение между собой 

элементов ГВВК; АРМ~--- автоматизированное рабочее место. 

      

      Эта эквивалентная схема представляет собой  

по\-сле\-до\-ва\-тель\-но-па\-рал\-лель\-ное соединение элементов ГВВК, участвующих 

в~решении <<информационных>> задач. 

      

      При последовательном соединении $n$ элементов системы результирующая 

вероятность работоспособности системы $P_p$ определяется~\cite{5-zac, 6-zac, 7-zac, 8-zac} 

как произведение 

вероятностей работоспособности входящих в~нее элементов: 



\noindent

$$

P_p=\prod\limits_{i=1}^n p_i\,. 

$$



\setcounter{figure}{1}

\begin{figure} %fig2

\vspace*{1pt}

 \begin{center}

 \mbox{%

 \epsfxsize=128mm 

 \epsfbox{zac-2.eps}

 }

 \end{center}

\vspace*{-9pt}

\Caption{Эквивалентная схема для расчета надежности ГВВК при решении 

<<информационных>> задач

}

\vspace*{-6pt}

\end{figure}



      При параллельном соединении $n$ элементов системы результирующая 

вероятность работоспособности системы~$P_p$ определяется через произведение 

вероятностей неработоспособности входящих в~нее элементов:



\noindent

$$

P_p=1- 

\prod\limits^n_{i=1} q_i\,,

$$

 где $q_i\hm=1\hm- p_i$~--- вероятность 

неработоспособности $i$-го элемента. 



\pagebreak

      

      В случае, когда $T_{\mathrm{о}}\hm\gg T_{\mathrm{в}}$, все элементы 

одинаковые ($\lambda_k\hm=\lambda$) и~отказы элементов независимы между собой, для 

оценки среднего времени безотказной работы  ($T_{\mathrm{op}}$) системы, состоящей из~$n$ параллельно 

соединенных элементов, можно воспользоваться выражением~\cite{6-zac}: 

      \begin{equation}

      T_{\mathrm{op}}=\fr{1}{\lambda} \sum\limits^n_{k=1} \fr{1}{k} = T_{\mathrm{о}}\left( 1+ 

\fr{1}{2} +\fr{1}{3}+\cdots +\fr{1}{n}\right)\,,

      \label{e1-zac}

\end{equation}

где $T_{\mathrm{о}}$~--- среднее время безотказной работы одного элемента; 

$T_{\mathrm{в}}$~--- среднее время восстановления работоспособности отказавшего 

элемента;  $\lambda\hm = 1/T_{\mathrm{о}}$~--- интенсивность отказов элемента 

в~каждый данный момент времени.



      Последовательность действий при вычислениях указана цифрами в~кружках.

      

      Сложнее с~эквивалентной схемой для расчета надежности ГВВК при решении 

задач <<подготовки данных>> (рис.~3). На рисунке пунктиром выделен фрагмент, 

который не может быть представлен в~виде по\-сле\-до\-ва\-тель\-но-па\-рал\-лель\-но\-го 

соединения элементов. Такие системы относятся к~системам с~произвольной 

структурой~[7, 8].



\begin{figure} %fig3

\vspace*{1pt}

 \begin{center}

 \mbox{%

 \epsfxsize=122.662mm 

 \epsfbox{zac-3.eps}

 }

 \end{center}

\vspace*{-9pt}

\Caption{Эквивалентная схема для расчета надежности ГВВК при решении задач 

<<подготовки данных>>}

\end{figure}



      Для расчета показателей надежности систем произвольной структуры разработан 

ряд точных и~приближенных методов~\cite{4-zac, 5-zac, 6-zac, 7-zac, 8-zac}. Наиболее 

простым из точных методов является метод прямого перебора состояний элементов 

системы. 

      

      Его суть состоит в~следующем: если определен критерий отказа системы, то все 

множество ее состояний можно разделить на два подмножества: подмножество состояний 

работоспособности~$F$ и~подмножество состояний отказа~$\Omega$. Чтобы найти 

вероятность безотказной работы системы с~произвольной структурой ($P_{\mathrm{ф}}$), 

необходимо сложить вероятности всех работоспособных состояний системы.

      

      Оценим надежность выделенного фрагмента указанным методом. 

      

      На рис.~4 для удобства вычислений представлен несколько преобразованный 

вариант эквивалентной схемы фрагмента, выделенного на рис.~3. На схеме элемент~1 

заменяет сервер~1 с~коэффициентом готовности~$p_1$, элементы~3--6 заменяют 

однотипные коммутаторы~3--6, обозначим их коэффициент готовности~$p_2$,\linebreak\vspace*{-12pt}



\pagebreak



\



\vspace*{-12pt}



\begin{floatingfigure}{65mm} %fig4

\vspace*{-8pt}

 \begin{center}

 \mbox{%

 \epsfxsize=61.571mm 

 \epsfbox{zac-4.eps}

 }

 \end{center}

\vspace*{-6pt}

\Caption{Эквивалентная схема фрагмента, выделенного на рис.~3}

\vspace*{14pt}

\end{floatingfigure}



\noindent 

а~элемент~2 заменяет сервер~3 и~коммутатор~1 с~результирующей на\-деж\-ностью двух 

последовательно соединенных элементов, обозначим его\linebreak

 коэффициент го\-тов\-ности~$p_3$.













      В табл.~1 показаны состояния \textit{работоспособности} выделенного фрагмента 

в~зависимости от состояния его элементов. Символами $x_1\mbox{--}x_6$ обозначены 

элементы~$1\mbox{--}6$, $x_i\hm = 1$ означает, что $i$-й элемент системы исправен, 

а~$x_i\hm=0$~--- что он неисправен. 



\begin{table}[b]\small %tabl1

      \begin{center}

      \Caption{Работоспособные состояния элементов системы }

      \vspace*{2ex}

      

      \tabcolsep=3pt

      \begin{tabular}{|c|c|c|c|c|c|c|c|c|c|c|c|c|c|c|c|}

      \hline  %\multicolumn{1}{|c|}{\raisebox{-6pt}[0pt][0pt]{

\multicolumn{1}{|c|}{\raisebox{-6pt}[0pt][0pt]

{\tabcolsep=0pt\begin{tabular}{c}Индекс\\ состо-\\ яния\end{tabular}}}&

\multicolumn{6}{c|}{Состояние элементов}&

\tabcolsep=0pt\begin{tabular}{c}Вероят-\\ ность\\ 

состо-\end{tabular}&

\multicolumn{1}{c|}{\raisebox{-6pt}[0pt][0pt]

{\tabcolsep=0pt\begin{tabular}{c}Индекс\\ состо-\\ яния\end{tabular}}}& 

\multicolumn{6}{c|}{Состояние элементов}

  &\tabcolsep=0pt\begin{tabular}{c}Вероят-\\ ность\\  состо-\end{tabular}\\

\cline{2-7}

\cline{10-15}

&$x_1$&$x_2$&$x_3$&$x_4$&$x_5$&$x_6$& яния&&$x_1$&$x_2$&$x_3$&$x_4$&$x_5$&$x_6$& яния\\

\hline

&&&&&&&&&&&&&&&\\[-9pt]

\hphantom{9}0&1&1&1&1&1&1&$p_1p_2p_3^4$&\hphantom{9}36&1&1&0&1&1&0&$p_1p_2q_3^2p_3^2$\\

&&&&&&&&&&&&&&&\\[-9pt]

\hphantom{9}1&0&1&1&1&1&1&$q_1p_2p_3^4$&\hphantom{9}45&1&1&1&0&0&1&$p_1p_2q_3^2p_3^2$\\

&&&&&&&&&&&&&&&\\[-9pt]

\hphantom{9}2&1&0&1&1&1&1&$p_1q_2p_3^4$&\hphantom{9}46&1&1&1&0&1&0&$p_1p_2q_3^2p_3^2$\\

&&&&&&&&&&&&&&&\\[-9pt]

\hphantom{9}3&1&1&0&1&1&1&$p_1p_2q_3p_3^3$&134&0&1&0&0&1&1&$q_1p_2q_3^2p_3^2$\\

&&&&&&&&&&&&&&&\\[-9pt]

\hphantom{9}4&1&1&1&0&1&1&$p_1p_2q_3p_3^3$&135&0&1&0&1&0&1&$q_1p_2q_3^2p_3^2$\\

&&&&&&&&&&&&&&&\\[-9pt]

\hphantom{9}5&1&1&1&1&0&1&$p_1p_2q_3p_3^3$&136&0&1&0&1&1&0&$q_1p_2q_3^2p_3^2$\\

&&&&&&&&&&&&&&&\\[-9pt]

\hphantom{9}6&1&1&1&1&1&0&$p_1p_2q_3p_3^3$&145&0&1&1&0&0&1&$q_1p_2q_3^2p_3^2$\\

&&&&&&&&&&&&&&&\\[-9pt]

13&0&1&0&1&1&1&$q_1p_2q_3p_3^3$&146&0&1&1&0&1&0&$q_1p_2q_3^2p_3^2$\\

&&&&&&&&&&&&&&&\\[-9pt]

14&0&1&1&0&1&1&$q_1p_2q_3p_3^3$&235&1&0&0& 1&0& 1&$p_1q_2q_3^2p_3^2$\\

&&&&&&&&&&&&&&&\\[-9pt]

15&0&1&1&1&0&1&$q_1p_2q_3p_3^3$&236&1&0&0&1 & 1&0&$p_1q_2q_3^2p_3^2$\\

&&&&&&&&&&&&&&&\\[-9pt]

16&0&1&1&1&1&0&$q_1p_2q_3p_3^3$&245&1&0& 1&0&0& 1&$p_1q_2q_3^2p_3^2$\\

&&&&&&&&&&&&&&&\\[-9pt]

23&1&0&0&1&1&1&$p_1q_2q_3p_3^3$&246&1&0& 1&0& 1&0&$p_1q_2q_3^2p_3^2$\\

&&&&&&&&&&&&&&&\\[-9pt]

24&1&0&1&0&1&1&$p_1q_2q_3p_3^3$&345&1&1&0&0&0&1&$p_1p_2q_3^3p_3$\\

&&&&&&&&&&&&&&&\\[-9pt]

25&1&0&1&1&0&1&$p_1q_2q_3p_3^3$&346&1&1&0&0&1&0&$p_1p_2q_3^3p_3$\\

&&&&&&&&&&&&&&&\\[-9pt]

26&1&0&1&1&1&0&$p_1q_2q_3p_3^3$&1345\hphantom{9}&0&1&0&0&0&1&$q_1p_2q_3^3p_3$\\

&&&&&&&&&&&&&&&\\[-9pt]

34&1&1&0&0&1&1&$p_1p_2q_3^2p_3^2$&1346\hphantom{9}&0&1&0&0&1&0&$q_1p_2q_3^3p_3$\\

&&&&&&&&&&&&&&&\\[-9pt]

35&1&1&0&1&0&1&$p_1p_2q_3^2p_3^2$&&&&&&&&\\

\hline

\end{tabular}

\end{center}

\end{table}

      

      Просуммировав вероятности исправных состояний системы, получим: 

      \begin{multline*}

P_{\mathrm{ф}} = p_1p_2p_3^4 + q_1p_2p_3^4 + p_1q_2p_3^4 + 4p_1p_2q_3p_3^3 + 

4q_1p_2q_3p_3^3 + 4p_1q_2q_3p_3^3 + {}\\[6pt]

{}+ 5p_1p_2q_3^2p_3^2 + 5q_1p_2q_3^2p_3^2 + 4p_1q_2q_3^2p_3^2 + 2p_1p_2q_3^3p_3 + 

2q_1p_2q_3^3p_3\,.

\end{multline*}



\setcounter{figure}{4}

 \begin{figure} %fig5

\vspace*{1pt}

 \begin{center}

 \mbox{%

 \epsfxsize=127.955mm 

 \epsfbox{zac-5.eps}

 }

 \end{center}

\vspace*{-9pt}

\Caption{Эквивалентная схема для расчета надежности ГВВК при решении 

<<расчетных>> задач}

\end{figure}



\pagebreak



\



\vspace*{-12pt}



\renewcommand{\figurename}{\protect\bf Таблица}

\setcounter{figure}{1}

\begin{floatingfigure}{75mm}  

\vspace*{-18pt}    

%      \begin{table}

{\small %tabl2

      \begin{center}

      \Caption{Исходные данные по на\-деж\-ности элементов ГВВК}

      \vspace*{2ex}

      

     \tabcolsep=9.5pt

      \begin{tabular}{|l|r|c|}

      \hline

\multicolumn{1}{|c|}{Устройство}&

\multicolumn{1}{c|}{$T_{\mathrm{о}}$, ч}&$1-K_{\mathrm{г}}$\\

\hline

&&\\[-9pt]

Серверы 1--3&220\,110&$1{,}7\cdot 10^{-7}$\\

Сервер Power 9&340\,200&$1{,}5\cdot 10^{-7}$\\

Коммутаторы 1, 2&525\,600&$3{,}4\cdot 10^{-6}$\\

Коммутаторы 3--6&613\,200&$2{,}9\cdot 10^{-6}$\\

Коммутаторы 7, 8&408\,800&$1{,}2\cdot 10^{-6}$\\

СХД&700\,000&$2{,}9\cdot 10^{-6}$\\

Средства ИБ&344\,200&$1{,}5\cdot10^{-6}$\\

СКС&20\,000&$2{,}5\cdot10^{-5}$\\

АРМ&12 000 &$4{,}2\cdot 10^{-5}$\\

\hline

\end{tabular}

\end{center}}

%\end{table}

%\vspace*{-12pt}

\end{floatingfigure}

      

\renewcommand{\figurename}{\protect\bf Рис.}

\renewcommand{\tablename}{\protect\bf Таблица}





      В остальном расчеты проводятся так же, как и~в~предыду\-щем случае (оценка 

на\-деж\-ности \mbox{ГВВК} при решении\linebreak

 <<информационных>> задач).\linebreak По\-сле\-до\-ва\-тель\-ность 

действий при вы\-чис\-ле\-ни\-ях указана циф\-ра\-ми в~круж\-ках.

      

      Эквивалентная схема для расчета на\-деж\-ности \mbox{ГВВК} при решении <<расчетных>> 

задач пред\-став\-ля\-ет собой по\-сле\-до\-ва\-тель\-но-па\-рал\-лель\-ное соединение 

элементов \mbox{ГВВК}, участ\-ву\-ющих в~решении <<рас\-чет\-ных>> задач (рис.~5). 

%  

 Расчет надежности этой схемы проводится аналогично расчету для случая решения 

<<информационных>> задач.

      

      Исходные данные для расчетов представлены в~табл.~2. 

      

     

\section{Оценка надежности и~анализ результатов расчетов}



      Используя предложенную методику и~представленные исходные данные, проведем 

расчеты надежности рассмотренных ФП. 

      

      В результате расчетов получены следующие результаты:

      \begin{itemize}

      \item  для ФП, решающей задачи <<информационного>> типа, вероятность отказа 

$G_{\mbox{{\tiny И}}}\hm = 7{,}2\cdot10^{-5}$;

      \item для ФП, решающей задачи типа <<подготовка данных>>, $G_{\mbox{\tiny ПД}}\hm 

= 7{,}5\cdot10^{-5}$; 

      \item для ФП, решающей задачи <<расчетного>> типа, $G_{\mbox{\tiny Р}}\hm = 

7{,}8\cdot10^{-5}$.

      \end{itemize}

      

      Кроме того, представляет интерес оценка надежности непосредственно 

<<вычислительной инфраструктуры>> ГВВК (без периферийных элементов), элементы 

которой в~основном резервируют друг друга (серверы~1 и~3, серверы~4 и~5, 

СХД~1 и~2, коммутаторы~3, 5 и~4, 6). Результаты расчетов вероятности отказа 

<<вычислительной инфраструктуры>> показали, что она изменяется в~пределах 

$G_{\mbox{\tiny ВИ}}\hm = 3{,}8\cdot 10^{-6}\ldots 2{,}3\cdot10^{-6}$ в~зависимости от 

конфигурации ее элементов при решении задач различных ФП. 

      

      В процессе решения научных и~научно-практических задач решаемая задача 

последовательно использует определенные ресурсы элементов ГВВК (АРМ, средства доступа, 

коммутаторы, серверы и~др.). Из теории надежности  

известно~\cite{5-zac, 6-zac, 7-zac, 8-zac}, что при последовательном соединении элементов 

результирующая надежность системы будет ниже надежности <<наихудшего>> с~точки 

зрения надежности элемента. В~данном случае надежность любой из рассмотренных ФП 

не может быть выше надежности АРМ.

      

      С другой стороны, ГВВК функционирует в~интересах большого числа 

подключенных к~нему пользователей, которых можно рассматривать как параллельно 

соединенные элементы, и~с~точки зрения <<равнопрочности>> надежность 

<<вычислительной инфраструктуры>> ГВВК должна быть не ниже надежности 

параллельно соединенных АРМ (т.\,е.\ если функционирует хотя бы один пользователь, то 

<<вычислительная инфраструктура>> также должна функционировать). Расчеты, 

проведенные с~использованием выражения~(1), показывают, что при подключении 

к~ГВВК от~50 до~100 пользователей результирующая вероятность отказа системы, 

состоящей из~$n$~параллельно соединенных АРМ, изменяется в~пределах 

$G_{\sum\mbox{\tiny АРМ}}\hm\approx9{,}3\cdot10^{-6}\ldots 8{,}1\cdot10^{-6}$, что выше вероятности 

отказа <<вычислительной инфраструктуры>>, и~с~большой долей уверенности можно 

утверждать, что равнопрочность этих двух составных частей ГВВК выполняется.

      

      Некоторое различие результатов оценки надежности для различных ФП 

объясняется тем, что при построении эквивалентных схем в~расчете на <<наихудший 

случай>> в~разных вариантах построения ФП задействуется разное число элементов 

системы (с~увеличением числа последовательно соединенных элементов результирующая 

надежность системы снижается).





      \section*{Выводы}

      

      В статье проведена оценка надежности 

ГВВК ФИЦ ИУ РАН при решении научных задач. По результатам 

исследований можно сделать следующие выводы:

      \begin{enumerate}[(1)]

      \item  как показали расчеты, среднее время наработки на отказ <<вычислительной 

инфраструктуры>> ГВВК ($\approx15$~лет) соизмеримо со средним сроком службы 

ГВВК. Это обеспечивается за счет использования в~структуре ГВВК высоконадежных 

элементов, а также за счет структурного резервирования наиболее важных элементов, 

участвующих в~решении задач, в~режиме <<горячего резервирования>>; 

      \item проведенная оценка вероятности отказа <<вычислительной 

инфраструктуры>> ГВВК показала, что она ниже вероятности отказа совокупности 

под\-клю\-ча\-емых к~ней АРМ, поэтому с~большой долей уверенности можно утверждать, что 

если функционирует хотя бы один АРМ пользователя, то <<вычислительная 

инфраструктура>> также будет функционировать;

      \item при необходимости обеспечить повышение надежности 

ФП наибольший прирост может быть получен при повышении надежности 

наименее надежных (<<наихудших>> с~точки зрения надежности) элементов (в данном 

случае АРМ и~СКС).

      \end{enumerate}

      

   {\small\frenchspacing

{%\baselineskip=10.8pt

\addcontentsline{toc}{section}{Литература}

\begin{thebibliography}{9}

\bibitem{1-zac}

\Au{Shan A.} Heterogeneous processing: A~strategy for augmenting Moore's law~/\!\!/

Linux~J., 2006. {\sf 

https://www.linuxjournal.com/article/8368}.

\bibitem{2-zac}

\Au{Соколов И.\,А., Зацаринный~А.\,А., Дивеев~А.\,И., Захаров~В.\,Н.,

Посыпкин~М.\,А., Абгарян~К.\,К.}

Гибридный высокопроизводительный вычислительный комплекс (\mbox{ГВВК}), 2018. {\sf 

http://www.frccsc.ru/hhpcc.}

\bibitem{3-zac}

ГОСТ 24.701-86. Единая система стандартов автоматизированных систем управления. 

Надежность автоматизированных систем управления. Основные положения. {\sf 

http://docs.cntd.ru/document/1200022035.}



\bibitem{5-zac} %4

\Au{Каштанов В.\,А., Медведев~А.\,И.} Теория надежности сложных систем.~--- М.: 

Физматлит, 2010. 608~с.

\bibitem{4-zac} %5

\Au{Зацаринный А.\,А., Гаранин А.\,И., Козлов~С.\,В.} На\-уч\-но-прак\-ти\-че\-ские 

аспекты обеспечения надежности  

ин\-фор\-ма\-ци\-он\-но-те\-ле\-ком\-му\-ни\-ка\-ци\-он\-ных сетей.~--- М.: ФИЦ ИУ РАН, 

2017. 246~с.

\bibitem{6-zac}

\Au{Гнеденко Б.\,В., Беляев~Ю.\,К., Соловьев~А.\,Д.} Математические методы в~теории 

надежности.~--- М.: Наука, 1965. 524~с.

\bibitem{7-zac}

\Au{Козлов Б.\,А., Ушаков~И.\,А.} Справочник по расчету надежности аппаратуры 

радиоэлектроники и~автоматики.~--- М.: Сов. радио, 1975. 462~с.

\bibitem{8-zac}

\Au{Беляев Ю.\,К., Богатырев~В.\,А., Болотин~В.\,В. и~др.} Надежность технических 

систем: Справочник~/ Под ред. И.\,А.~Ушакова.~--- М.: Радио и~связь, 1985. 608~с.

        \end{thebibliography}

} }





\vspace*{-3pt}



\hfill{\small\textit{Поступила в~редакцию 15.03.19}}









%\vspace*{8pt}



%\hrule



%\vspace*{2pt}



\newpage



\vspace*{-28pt}



%\hrule



%\vspace*{9pt}







\def\tit{EVALUATION OF~RELIABILITY OF~THE~HYBRID HIGH-PERFORMANCE 

COMPUTING COMPLEX\\ IN~SOLUTION OF~SCIENTIFIC PROBLEMS}



\def\titkol{Evaluation of reliability of the hybrid high-performance 

computing complex} % in solution of scientific problems}





\def\aut{A.\,A.~Zatsarinny$^1$, A.\,I.~Garanin$^1$, V.\,A.~Kondrashev$^1$, K.\,I.~Volovich$^1$, 

and~S.\,I.~Malkovsky$^2$}

\def\autkol{A.\,A.~Zatsarinny, A.\,I.~Garanin, V.\,A.~Kondrashev, et al.}

%K.\,I.~Volovich$^1$,  and~S.\,I.~Malkovsky$^2$}







\titel{\tit}{\aut}{\autkol}{\titkol}



\vspace*{-9pt}



 \noindent

$^1$Institute of Informatics Problems, Federal Research Center ``Computer 

Science\linebreak

$\hphantom{^1}$and Control'' of the Russian Academy of Sciences, 44-2~Vavilov Str., 

Moscow\linebreak

$\hphantom{^1}$119133, Russian Federation



\noindent

$^2$Computing Center, Far Eastern Branch of the Russian Academy of Sciences,\linebreak 

$\hphantom{^1}$65~Kim~U~Chen Str., Khabarovsk 680000, Russian Federation



 \noindent



\def\leftfootline{\small{\textbf{\thepage}

\hfill Sistemy i Sredstva Informatiki~---

Systems and Means of Informatics 2019 vol~29 no\,2}

}%

 \def\rightfootline{\small{Sistemy i Sredstva Informatiki~---

 Systems and Means of Informatics   2019   vol~29   no\,2

\hfill \textbf{\thepage}}}







\Abste{The necessity of using hybrid solutions in creation of high-performance computing 

systems is substantiated. A brief description of the hybrid high-performance computing complex 

(HHPCC) of FRC CSC RAS is given and its enlarged block diagram is presented. The paper offers 

the methodical approach to estimation of reliability of HHPCC on the basis of which calculations of 

reliability of the allocated functional subsystems are carried out. Separately, reliability of the 

``computing infrastructure'' of HHPCC (without peripheral elements) was evaluated. 

Recommendations for improving reliability of functional subsystems are given.}



\KWE{hybrid high-performance computing system; reliability; functional subsystems; failure; 

equivalent circuit for reliability calculation}



\DOI{10.14357\!/\!08696527190212} 



\vspace*{-6pt}



\Ack



\vspace*{-2pt}

\noindent

The work was partly supported by 

the Russian Foundation for Basic Research (projects 18-29-03100

and 18-29-03196). 





\renewcommand{\bibname}{\protect\rmfamily References}

%\renewcommand{\bibname}{\large\protect\rm References}



\vspace*{-2pt}



{\small\frenchspacing

{%\baselineskip=10.8pt

\addcontentsline{toc}{section}{References}

\begin{thebibliography}{9}



\vspace*{-2pt}



\bibitem{1-zac-1}

\Aue{Shan, A.} 2006. Heterogeneous processing: A~strategy for augmenting Moore's law. 

\textit{Linux~J.}

Available at: {\sf https://www.linuxjournal.com/article/8368} (accessed April~4, 2019).

\bibitem{2-zac-1}

\Aue{Sokolov, I.\,A., A.\,A.~Zatsarinnyy, A.\,I.~Diveev, V.\,N.~Zakharov, M.\,A.~Posypkin, 

and K.\,K.~Abgarjan.} 2018. Gibridnyy 

vysokoproizvoditel'nyy vychislitel'nyy kompleks (GVVK) 

[Hybrid high-performance computing system (HHPCC)]. Available at: 

{\sf http://www.frccsc.ru/hhpcc} 

(accessed April~4, 2019).

\bibitem{3-zac-1}

GOST 24.701-86. 1986. Edinaya sistema standartov avtomatizirovannykh sistem upravleniya. 

Nadyozhnost' avtomatizirovannykh sistem upravleniya. Osnovnye polozheniya 

[State Standard ``An 

uniform system of standards for automated control systems. Reliability of automated control 

systems. Fundamentals'']. Available at: {\sf http://docs.cntd.ru/document/1200022035} (accessed 

April~4, 2019).

\bibitem{5-zac-1}

\Aue{Kashtanov, V.\,A., and A.\,I.~Medvedev.} 2010. \textit{Teoriya nadezhnosti slozhnykh 

sistem} [The reliability theory for complex systems]. Moscow: Fizmatlit. 608~p.



\bibitem{4-zac-1} %5

\Aue{Zatsarinnyy, A.\,A., A.\,I.~Garanin, and S.\,V.~Kozlov.} 2017.  

\textit{Nauchno-prakticheskie aspecty obespecheniya nadezhnosti  

informatsionno-telekommunikatsionnykh setey} [Scientific and practical aspects of ensuring 

reliability of information and telecommunication networks]. Moscow: FIC IU RAN. 246~p.

\bibitem{6-zac-1}

\Aue{Gnedenko, B.\,V., Y.\,K.~Belyaev, and A.\,D.~Soloviev.} 1965. \textit{Matematicheskie 

metody v~teorii nadezhnosti} [Mathematical methods in the theory of reliability]. Moscow: 

Nauka.  524~p.

\bibitem{7-zac-1}

\Aue{Kozlov, B.\,A., and I.\,A.~Ushakov.} 1975. \textit{Spravochnik po raschetu nadezhnosti 

apparatury radioelektroniki i~avtomatiki} [Handbook on calculation of 

reliability of radioelectronics and automation equipment]. Moscow: Sovetskoe Radio. 462~p.

\bibitem{8-zac-1}

\Aue{Belyaev, Yu.\,K., V.\,A.~Bogatyrev, V.\,V.~Bolotin, \textit{et al.}} 1985. Nadezhnost' 

tekh\-ni\-che\-skikh sistem: Spravochnik [Reliability of technical systems: Handbook]. Ed. 

I.\,A.~Ushakov. Moscow: Radio i~svyaz'. 608~p.

\end{thebibliography}

} }





\vspace*{-3pt}



\hfill{\small\textit{Received March 15, 2019}} 



\Contr



\noindent

\textbf{Zatsarinny Alexander A.} (b.\ 1951)~--- Doctor of Science in technology, 

professor, Deputy Director, Institute of Informatics Problems, Federal Research 

Center ``Computer Science and Control'' of the Russian Academy of Sciences, 44-

2~Vavilov Str., Moscow 119133, 

Russian Federation; \mbox{AZatsarinny@frccsc.ru}



\vspace*{3pt}



\noindent

\textbf{Garanin Alexander I.} (b.\ 1951)~--- Candidate of Science (PhD) in technology, 

senior scientist, Institute of Informatics Problems, Federal Research 

Center ``Computer Science and Control'' of the Russian Academy of Sciences, 

 44-2~Vavilov Str., Moscow 119133, Russian Federation; \mbox{algaranin@mail.ru}



\vspace*{3pt}



\noindent

\textbf{Kondrashev Vadim A.} (b.\ 1963)~--- senior scientist, Institute of Informatics 

Problems, Federal Research Center ``Computer Science and Control'' of the Russian Academy of 

Sciences, 44-2~Vavilov Str., Moscow 119133, Russian Federation; \mbox{vd@ipi.ac.ru}



\vspace*{3pt}



\noindent

\textbf{Volovich Konstantin I.} (b.\ 1970)~--- senior scientist, Institute of Informatics 

Problems, Federal Research Center ``Computer Science and Control'' of the Russian 

Academy of Sciences, 44-2~Vavilov Str., Moscow 119133, Russian Federation;  

\mbox{kv@ipi.ac.ru}



\vspace*{3pt}



\noindent

\textbf{Malkovsky Sergey I.} (b.\ 1983)~--- scientist, Computing Center, Far Eastern Branch of 

the Russian Academy of Sciences, 65~Kim~U~Chen Str., Khabarovsk 680000, Russian 

Federation; \mbox{sergey.malkovsky@gmail.com} 





\label{end\stat}





\renewcommand{\bibname}{\protect\rm Литература}

